\documentclass{article}

\input{~/University/hwpreamble.tex}

\begin{document}

\title{Estadística: Tema 4}
\author{Victoria Eugenia Torroja Rubio}
\date{\today}

\maketitle

\begin{ej}
Método de los momentos: Distribución normal. Sea $\displaystyle \left(X_{1}, \ldots, X_{n}\right) $ una m.a.s. de $\displaystyle X \sim N\left(\mu, \sigma^{2}\right) $. Tenemos que los momentos poblacionales son
\[\alpha_{1}=E\left(X\right)=\mu \quad \text{y} \quad \alpha_{2} = E\left(X^{2}\right)=\sigma^{2} + \mu^{2} .\]
Por otro lado, los momentos muestrales son
\[a_{1}=\frac{1}{n}\sum^{n}_{i = 1}X_{i}=\overline{X} \quad \text{y} \quad a_{2}=\frac{1}{n}\sum^{n}_{i = 1}X_{i}^{2} .\]
Tenemos que el estimador para la media será
\[\hat{\mu}=\overline{X} \quad \text{y} \quad \sigma^{2} + \mu^{2} = a_{2} \Rightarrow \hat{\sigma }_{2}=a_{2}-\mu^{2} = \frac{1}{n}\sum^{n}_{i =1}X_{i}^{2}-\overline{X}^{2} = S^{2} .\]
\end{ej}
\begin{ej}

\end{ej}


\end{document}
